\begin{tikzpicture}[font=\small]
\node[anchor=north west, inner sep=0] at (0,0) {%
\begin{tabular}{llll}
\toprule
\texttt{i\_op} & name & result \texttt{o\_y} & flags that can move\\
\midrule
\texttt{000} & \texttt{ADD} & $a+b$ & Z N \textbf{C V}\\
\texttt{001} & \texttt{SUB} & $a-b$ & Z N \textbf{C V}\\
\texttt{010} & \texttt{AND} & $a\mathbin{\&}b$ & Z N\\
\texttt{011} & \texttt{OR}  & $a\mathbin{|}b$ & Z N\\
\texttt{100} & \texttt{XOR} & $a\mathbin{\hat{}}b$ & Z N\\
\texttt{101} & \texttt{SHL} & $a\ll1$ & Z N \textbf{C}\\
\texttt{110} & \texttt{SHR} & $a\gg1$ & Z N \textbf{C}\\
\texttt{111} & \texttt{CMP} & $a-b$ & Z N \textbf{C V}\\
\bottomrule
\end{tabular}};
\node[anchor=north west, text width=6.2cm, font=\scriptsize, text=slate]
  at (0,-5.6) {The encoding is a contract, not a preference: the ch16
  instruction set drives arithmetic through these three bits, so the
  table is written in stone for the rest of the book. \texttt{CMP}
  shares the subtracter's datapath but the difference it computes is
  thrown away --- the flags are the verdict.};
\end{tikzpicture}
